%% Đề 17 -- de-kiem-tra-cuoi-hoc-ki-2-toan-9-nam-2021-2022-so-gddt-lang-son.pdf
\deheader{ĐỀ SỐ 17}%
  {Đề kiểm tra Học kỳ — Năm học 2021--2022}%
  {Sở GD\&ĐT Lạng Sơn | Chỉ phần Đại số}
\addcontentsline{toc}{section}{Đề 17 — Đề kiểm tra Học kỳ — Năm học 2021--2022}

\begin{de_cau}{Câu 1}{2,0 điểm}
\begin{enumerate}[label=\alph*)]
\item Giải phương trình: $x^2 - 3x - 10 = 0$.
\item Cho $A=\left(\dfrac{1}{\sqrt{x}+2}+\dfrac{1}{\sqrt{x}-2}\right):\dfrac{2}{\sqrt{x}-2}$, $x\geq 0, x\neq 4$.
Rút gọn $A$.
\end{enumerate}
\end{de_cau}

\begin{loi_giai_cau}
\textbf{a)} $\Delta=9+40=49$. $x=\dfrac{3\pm 7}{2}$: $x=5$ hoặc $x=-2$.

\textbf{b)}
\[A=\frac{\sqrt{x}-2+\sqrt{x}+2}{(\sqrt{x}+2)(\sqrt{x}-2)}\cdot\frac{\sqrt{x}-2}{2}
=\frac{2\sqrt{x}}{x-4}\cdot\frac{\sqrt{x}-2}{2}=\frac{\sqrt{x}}{\sqrt{x}+2}.\]
\end{loi_giai_cau}

\begin{de_cau}{Câu 2}{2,0 điểm}
\begin{enumerate}[label=\alph*)]
\item Tìm tọa độ giao điểm của $(P)\colon y=x^2$ và $(d)\colon y=-x+2$.
\item Cho hàm số bậc nhất $y=(m-2)x+3$ ($m\neq 2$). Tìm $m$ để đồ thị đi qua $A(1;5)$.
\end{enumerate}
\end{de_cau}

\begin{loi_giai_cau}
\textbf{a)} $x^2=-x+2\Rightarrow x^2+x-2=0\Rightarrow(x+2)(x-1)=0$.
$x=-2,y=4$: điểm $A(-2;4)$; $x=1,y=1$: điểm $B(1;1)$.

\textbf{b)} Đồ thị qua $A(1;5)$: $5=(m-2)\cdot 1+3=m+1\Rightarrow m=4$.
\end{loi_giai_cau}

\begin{de_cau}{Câu 3}{1,5 điểm}
Giải bài toán bằng cách lập hệ phương trình:\\
Hai loại hàng có tổng 60 kg, trị giá 1 200 000 đồng. Biết giá loại 1 là 25 000 đồng/kg và loại 2 là
15 000 đồng/kg. Tìm số kg mỗi loại.
\end{de_cau}

\begin{loi_giai_cau}
Gọi số kg loại 1 là $x$, loại 2 là $y$ ($x,y>0$).

$\begin{cases}x+y=60\\25000x+15000y=1200000\end{cases}\Rightarrow\begin{cases}x+y=60\\5x+3y=240\end{cases}$.

Từ (1): $y=60-x$. Thế vào: $5x+3(60-x)=240\Rightarrow 2x=60\Rightarrow x=30, y=30$.

Mỗi loại \textbf{30 kg}.
\end{loi_giai_cau}

\begin{de_cau}{Câu 4}{1,0 điểm}
Cho phương trình $x^2-2mx+m+2=0$ (tham số $m$). Tìm $m$ để phương trình có hai nghiệm $x_1,x_2$
thỏa mãn $\dfrac{1}{x_1}+\dfrac{1}{x_2}=\dfrac{1}{2}$.
\end{de_cau}

\begin{loi_giai_cau}
Để có hai nghiệm phân biệt: $\Delta'=m^2-(m+2)>0\Rightarrow m^2-m-2>0\Rightarrow(m-2)(m+1)>0$.

$m>2$ hoặc $m<-1$.

Vi-ét: $x_1+x_2=2m$, $x_1x_2=m+2\neq 0$ (tức $m\neq -2$).

$\dfrac{1}{x_1}+\dfrac{1}{x_2}=\dfrac{x_1+x_2}{x_1x_2}=\dfrac{2m}{m+2}=\dfrac{1}{2}\Rightarrow 4m=m+2\Rightarrow m=\dfrac{2}{3}$.

Kiểm tra: $m=\frac{2}{3}$ không thỏa $m>2$ hay $m<-1$. Vậy không có $m$ thỏa mãn.
\end{loi_giai_cau}
